\section{Model and computable information}
\label{sec:model}
Let $Q\subset\mathbb R^d$ be a nonempty compact convex set with
$\operatorname{diam}Q\le D$, where $D>0$. Let $f$ be convex on $Q$
and twice continuously differentiable on a neighborhood of $Q$. Assume
\begin{equation}\label{eq:gradient_smoothness}
 \|\nabla^2f(x)\|\le L_0+L_1\|\nabla f(x)\|,\qquad x\in Q,
 \quad L_0,L_1\ge0.
\end{equation}
Vector norms are Euclidean and matrix norms are spectral. A minimizer
$x^*\in Q$ exists. It is used to state the accuracy guarantee, not as an
input to the method. We write $f(x)-f(x^*)$ and $\|\nabla f(x)\|$
explicitly, without introducing names for these two expressions.

The available operations are objective evaluation, gradient evaluation,
and linear minimization over all of $Q$. The latter returns
\[
 s\in\arg\min_{v\in Q}\langle\nabla f(x),v\rangle.
\]
The full Frank-Wolfe gap is the computed quantity
\begin{equation}\label{eq:gap}
 G(x)=\max_{v\in Q}\langle\nabla f(x),x-v\rangle
     =\langle\nabla f(x),x-s\rangle.
\end{equation}
Since $x\in Q$, $G(x)\ge0$. Convexity gives
$f(x)-f(x^*)\le G(x)$, so $f(x)-G(x)$ is a valid lower bound on the
unknown minimum value \cite{jaggi}. In particular, the stopping test
does not evaluate that minimum.

At a constrained minimizer the correct first-order condition is
\[
 \langle\nabla f(x^*),v-x^*\rangle\ge0\quad(v\in Q),
\]
which is equivalent to $G(x^*)=0$. It does not require
$\nabla f(x^*)=0$. For example, $f(x)=Ax+x^2/2$ on $[0,1]$, $A>0$,
has its minimum at zero, derivative $A$ there, and Hessian equal to one.
A large gradient and large curvature are therefore different phenomena.

\begin{lemma}[A finite curvature bound from the initial gradient]
\label{lem:global_curvature}
The smoothness condition implies
\begin{equation}\label{eq:global_curvature}
 \|\nabla^2f(z)\|\le K_Q\quad(z\in Q),\qquad
 K_Q=(L_0+L_1\|\nabla f(x_0)\|)e^{L_1D}.
\end{equation}
Any sharper valid uniform bound can replace $K_Q$ in the results below.
This constant is used to express work bounds; the algorithm does not
need to know it.
\end{lemma}
\begin{proof}
For a unit direction $e$ along a segment from $x_0$ in $Q$, the
fundamental theorem of calculus gives
\[
 \|\nabla f(x_0+te)\|\le\|\nabla f(x_0)\|
 +\int_0^t\{L_0+L_1\|\nabla f(x_0+se)\|\}\,ds.
\]
For $L_1>0$, Gronwall's inequality bounds
$L_0+L_1\|\nabla f(x_0+te)\|$ by
$(L_0+L_1\|\nabla f(x_0)\|)e^{L_1t}$. Since $t\le D$, this proves
the assertion. If $L_1=0$, the original condition directly gives the
bound $L_0$.
\end{proof}

The smoothness condition contains no minimum value. The algorithm below
also uses neither a regularity constant centered at an unknown optimizer
nor a positive-distance assumption on that optimizer. Convexity is the
property that makes the full gap a global error certificate.
